\documentclass{article}
\usepackage[backend=biber,natbib=true,style=alphabetic,maxbibnames=50]{biblatex}
\addbibresource{/home/nqbh/reference/bib.bib}
\usepackage[utf8]{vietnam}
\usepackage{tocloft}
\renewcommand{\cftsecleader}{\cftdotfill{\cftdotsep}}
\usepackage[colorlinks=true,linkcolor=blue,urlcolor=red,citecolor=magenta]{hyperref}
\usepackage{amsmath,amssymb,amsthm,enumitem,float,graphicx,mathtools,tikz}
\usetikzlibrary{angles,calc,intersections,matrix,patterns,quotes,shadings}
\allowdisplaybreaks
\newtheorem{assumption}{Assumption}
\newtheorem{baitoan}{Bài toán}
\newtheorem{cauhoi}{Câu hỏi}
\newtheorem{conjecture}{Conjecture}
\newtheorem{corollary}{Corollary}
\newtheorem{dangtoan}{Dạng toán}
\newtheorem{definition}{Definition}
\newtheorem{dinhly}{Định lý}
\newtheorem{dinhnghia}{Định nghĩa}
\newtheorem{example}{Example}
\newtheorem{ghichu}{Ghi chú}
\newtheorem{hequa}{Hệ quả}
\newtheorem{hypothesis}{Hypothesis}
\newtheorem{lemma}{Lemma}
\newtheorem{luuy}{Lưu ý}
\newtheorem{nhanxet}{Nhận xét}
\newtheorem{notation}{Notation}
\newtheorem{note}{Note}
\newtheorem{principle}{Principle}
\newtheorem{problem}{Problem}
\newtheorem{proposition}{Proposition}
\newtheorem{question}{Question}
\newtheorem{remark}{Remark}
\newtheorem{theorem}{Theorem}
\newtheorem{vidu}{Ví dụ}
\usepackage[margin=2cm]{geometry}
\def\labelitemii{$\circ$}
\DeclareRobustCommand{\divby}{%
	\mathrel{\vbox{\baselineskip.65ex\lineskiplimit0pt\hbox{.}\hbox{.}\hbox{.}}}%
}
\setlist[itemize]{leftmargin=*}
\setlist[enumerate]{leftmargin=*}

\title{Đề Thi Cuối Kỳ Tổ Hợp {\it\&} Lý Thuyết Đồ Thị Hè 2026\\Final Exam: Combinatorics {\it\&} Graph Theory Summer 2026}
\author{Giảng viên ra đề: Nguyễn Quản Bá Hồng}
\date{\today}

\begin{document}
\maketitle
\begin{center}\large
    \textbf{\textsf{Yêu cầu Bắt Buộc}}
\end{center}

\begin{enumerate}
    \item Được phép sử dụng tài liệu giấy (sách, báo, tập{\tt/}nháp viết tay, khăn giấy, giấy súc, etc.) không giới hạn.
    \item Cấm sử dụng thiết bị điện tử (trừ máy tính bỏ túi không có tích hợp AIs), AI tools (Gemini, ChatGPT, Grok, etc.). Nếu phát hiện 0 điểm ngay lần đầu tiên \& thu bài -- không có bất cứ cảnh cáo nào.
    \item Nếu làm không được, viết định nghĩa sẽ được +2 điểm. Trong mỗi bài, nếu không làm được ý trước thì vẫn có thể sử dụng kết quả các ý trước để làm ý sau của bài toán.
    \item Thời gian thi: 2.5 hours. Có thể nộp bài sớm nếu tự sinh viên muốn hoặc bị giảng viên bắt. Nộp bài sớm vẫn phải ngồi lại phòng thi, không được về sớm, để chiêm nghiệm về thái độ học tập, làm việc của bản thân \& đặc biệt về cuộc đời.
    \item Nếu phát hiện đề sai rồi sửa lại được đề \& làm đúng thì sẽ được cộng thêm nhiều điểm.
    \item Đề thi tổng cộng $\Sigma = 100$ điểm.
\end{enumerate}

\begin{center}\large
    \textbf{\textsf{Đề Thi Chính Thức}}
\end{center}

\begin{baitoan}[Basic graphs -- Các đồ thị cơ bản]
	(a) \emph{(5 điểm)} Cho ví dụ \& vẽ đồ thị đơn, đa đồ thị, đồ thị tổng quát có $n$ đỉnh với mỗi $n\in[5]$. Viết cụ thể tập đỉnh $V$ \& tập $E$ cho từng ví dụ.
	\item(b) \emph{(10 điểm)} Với đồ thị đơn $n\in\mathbb{N}^\star$ đỉnh tổng quát, tính số cạnh lớn nhất có thể. Viết $|V|,|E|$, $\{\deg v;v\in V\}$ trong trường hợp cực đại đó. Giải thích tại sao không thể xác định được với đa đồ thị hay đồ thị tổng quát?
\end{baitoan}

\begin{baitoan}[Some special graphs -- Vài đồ thị đặc biệt]
	\emph{(15 điểm)} Vẽ, viết ra $V,E,|V|,|E|$ cho đồ thị: (a) $P_{n-1}$. (b) Cycle $C_n$. (c) Complete graph $K_n$. (d) Complete bipartite graph $K_{m,n}$.
\end{baitoan}

\begin{baitoan}[Integer partition -- Phân hoạch số nguyên]
	(a) \emph{(5 đ)} Viết định nghĩa phân hoạch số nguyên của $n$, biểu đồ Ferrers, \& cho ví dụ. (b) \emph{(5 đ)} Viết định nghĩa phân hoạch liên hợp, biểu đồ Ferrer tương ứng, \& cho ví dụ.	
\end{baitoan}

\begin{baitoan}[Pentagon number -- Số ngũ giác]
	(a) \emph{(5 đ)} Viết định nghĩa \& viết công thức toán của $|s(\lambda)|$ \& $|l(\lambda)|$.
	\item(b) \emph{(5 đ)} Phát biểu định lý Euler cho số ngũ giác.
	\item(c) \emph{(10 điểm)} Tính giá trị của hàm số $f(n)\coloneqq|{\cal E}(n)| - |{\cal O}(n)|$ với $n\in[10]$.
\end{baitoan}

\begin{baitoan}[Regular graph -- Đồ thị chính quy]
	\emph{(10 đ)} Vẽ đồ thị $d$-chính quy có $7$ đỉnh với mỗi $d\in[6]$. Nếu không thể vẽ được về mặt toán học, chứng minh tại sao không vẽ được.
\end{baitoan}

\begin{baitoan}[Euler's theorem{\tt/}handshaking lemma {\it\&} Havel--Hakimi algorithm -- Định lý Euler/bổ đề bắt tay {\it\&} thuật toán Havel--Hakimi]
	{\rm($\Sigma = 15$ điểm)}
	\item(a) {\rm(5 điểm)} Phát biểu định lý Euler \& định lý về thuật toán Havel--Hakimi. 2 định lý này áp dụng cho các loại đồ thị nào?
	\item(b) {\rm(5 điểm) \cite[P10.1.13, p. 368]{Shahriari2022}} Chứng minh 1 dãy $d_1,d_2,\ldots,d_n$ là 1 graphic sequence khi \& chỉ khi dãy $n - d_n - 1,\ldots,n - d_2 - 1,n - d_1 - 1$ là graphic; áp dụng kiểm tra
	\begin{equation*}
		14,13,12,11,10,9,9,9,9,9,9,9,9,8,8,8,7,6,5,4,3,2
	\end{equation*}
	có phải là graphic sequence không.
	\item(c) {\rm(5 điểm)} Sử dụng thuật toán Havel--Hakimi để kiểm tra dãy
	\begin{equation*}
		14,13,12,11,10,9,9,9,9,9,9,9,9,8,8,8,7,6,5,4,3,2
	\end{equation*}
	có phải là graphic sequence hay không.
\end{baitoan}

\begin{baitoan}[Paths, trees {\it\&} degree sequences -- Đường đi, cây, {\it\&} dãy bậc]
	{\rm($\Sigma = 15$ điểm)} Trong khoa học máy tính, 1 đồ thị cây (tree) $T = (V, E)$ là 1 đồ thị đơn vô hướng, liên thông (connected) \& không chứa chu trình (cycle). Đã biết, mọi đồ thị cây có $n$ đỉnh thì luôn có đúng $|E| = n - 1$ cạnh.	
	\item(a) {\rm(5 điểm)} Viết định nghĩa của đường đi (path) \& chu trình (cycle) trong đồ thị.
	\item(b) {\rm(5 điểm)} Kế thừa Định lý bắt tay (Handshaking lemma) ở Bài 6, hãy chứng minh: Mọi đồ thị cây có $n \ge 2$ đỉnh luôn tồn tại ít nhất $2$ đỉnh treo (đỉnh lá -- leaf, tức là đỉnh có bậc đúng bằng $1$).
	\item(c) {\rm(5 điểm)} \emph{(Tư duy thuật toán)} Chứng minh: 1 dãy gồm $n$ số nguyên dương $d_1\ge d_2\ge\ldots\ge d_n\ge1$ ($n \ge 2$) là dãy bậc (degree sequence) của 1 đồ thị cây khi \& chỉ khi:
	\begin{equation*}
		\sum_{i = 1}^n d_i = 2n - 2.
	\end{equation*}
	{\sf Hint.} Chiều đảo $\Leftarrow$: Sử dụng quy nạp theo số đỉnh $n$. Nếu tổng các bậc bằng $2n - 2$ \& tất cả $d_i \ge 1$, chắc chắn phải có ít nhất 1 phần tử bằng 1. Thử ``ngắt'' đỉnh lá này đi \& nối vào 1 đồ thị cây $n-1$ đỉnh).
\end{baitoan}

%------------------------------------------------------------------------------%

\printbibliography[heading=bibintoc]

\end{document}